\documentclass[a4paper]{article}

\usepackage{amsmath,amssymb,amsthm}
\usepackage{fullpage}
\usepackage{hyperref}

\newtheorem{theorem}{Theorem}[section]
\newtheorem{lemma}[theorem]{Lemma}
\newtheorem{corollary}[theorem]{Corollary}
\newtheorem{remark}[theorem]{Remark}
\newtheorem{problem}[theorem]{Unresolved assertion}

\newcommand{\F}{\mathbb F}
\newcommand{\Z}{\mathbb Z}
\newcommand{\ord}{\operatorname{ord}}
\DeclareMathOperator{\Tr}{Tr}
\DeclareMathOperator{\Nm}{N}

\setlength{\parindent}{0pt}
\setlength{\parskip}{0.7em}

\title{Two Trace--Norm Theorems for Prime-Power\\
Lights Out Stabilization}
\author{}
\date{September 5, 2026}

\begin{document}
\maketitle

\begin{abstract}
For every prime $p>5$, we prove $d(p^k-1)=0$ for all $k\geq1$
whenever either $3\mid\ord_p(2)$ or $v_2(\ord_p(2))=2$.
These results use trace--norm identities in cubic norm-one groups and
in the two coprime factors of $2^{2r}+1$ for odd $r$.
They apply to arbitrary Wieferich depth and give analytic proofs for
both known Wieferich primes, $1093$ and $3511$, without their endpoint
searches.  We also derive an exact nullity-increment formula which
separates new--new from new--old root collisions.

The identity $d(p^k-1)=d(p-1)$ for \emph{every} prime is not proved
here, and no counterexample is found.  The remaining case is an
arbitrary Wieferich prime with
$3\nmid\ord_p(2)$ and $v_2(\ord_p(2))\ne2$.
\end{abstract}

\section{Definitions and relation to the existing work}

All polynomial calculations, unless otherwise specified, take place in
$\F_2[X]$.  Let
\[
 F_0=0,\qquad F_1=1,\qquad F_{n+1}=XF_n+F_{n-1}.
\]
For the bounded $n\times n$ Lights Out grid, the existing writeups use
Sutner's formula
\begin{equation}\label{eq:nullity}
 d(n)=\deg\gcd(F_{n+1}(X),F_{n+1}(X+1)).
\end{equation}
In particular $d(0)=0$.  Throughout, the exponent $k$ in the proposed
identity is positive.
We write $\mu_N=\{z\in\overline{\F}_2:z^N=1\}$,
$\ord_m(2)$ for the multiplicative order of $2$ modulo an odd integer
$m>1$, and $v_\ell$ for the usual $\ell$-adic valuation.

The principal sources within this repository are
\cite{finite-fields,plateau}.  They already prove stabilization above
$c_p=v_p(2^{\ord_p(2)}-1)$, give the endpoint equivalence, and analyze
many unsuccessful approaches to the plateau.
Section~\ref{sec:norms} supplies two additional analytic families,
including both known Wieferich primes.
Section~\ref{sec:increment} gives a quantitative refinement of the
endpoint gcd criterion.  Section~\ref{sec:descent} includes the known
descent argument to make the precise remaining gap independently
checkable.

For odd $N\geq1$, define
\[
 H_N=F_{(N-1)/2}+F_{(N+1)/2},\qquad
 G_N=\gcd(H_N(X),H_N(X+1)).
\]
The Fibonacci identities give $F_N=H_N^2$, and consequently
\begin{equation}\label{eq:half-nullity}
 d(N-1)=2\deg G_N.
\end{equation}
Here $H_1=G_1=1$.

\begin{lemma}\label{lem:roots}
For odd $N$, the polynomial $H_N$ is monic and square-free, with root set
\[
 T_N=\{u+u^{-1}:u^N=1,\ u\ne1\}\subset\overline{\F}_2.
\]
If $M\mid N$ is odd, then $H_M\mid H_N$.
\end{lemma}

\begin{proof}
The identity
\[
 (u+u^{-1})F_N(u+u^{-1})=u^N+u^{-N}
\]
shows that every displayed value is a root.  For $u\ne1$ of odd order,
$u+u^{-1}\ne0$.  Two distinct nonzero elements give the same value precisely
when they are mutual inverses: they are the roots of
$Z^2+(u+u^{-1})Z+1$.
Thus there are $(N-1)/2$ distinct displayed roots, exactly the degree
of the monic polynomial $H_N$.  This proves square-freeness and the
description of its roots.  The inclusion $T_M\subseteq T_N$ proves
divisibility.
\end{proof}

\section{Two trace--norm exclusions}\label{sec:norms}

By Lemma~\ref{lem:roots}, a common root is equivalent to the existence
of $u,v\ne1$ in $\mu_N$ satisfying
\begin{equation}\label{eq:relation}
 1+u+u^{-1}+v+v^{-1}=0.
\end{equation}
The exclusions below treat all pairs, including pairs of different
orders and different Frobenius orbits.

\subsection{The cubic norm-one group}

\begin{theorem}[Cubic exclusion]\label{thm:cubic}
Let $q$ be a power of two.  Suppose $u,v\in\F_{q^3}^{\times}$ have
norm one to $\F_q$ and satisfy \eqref{eq:relation}.
Then either $u=1$ and $v$ has order $3$, or $v=1$ and $u$ has
order $3$.  In particular,
\begin{equation}\label{eq:cubic-nullity}
 d(q^2+q)=0.
\end{equation}
\end{theorem}

\begin{proof}
Write $\Tr$ and $\Nm$ for the trace and norm from $\F_{q^3}$ to
$\F_q$, and set $a=u+u^{-1}$.  Since $\Nm(u)=1$,
\[
 a=\frac{(u+1)^2}{u},\qquad
 \Nm(a)=\Nm(u+1)^2.
\]
Expanding the product defining $\Nm(u+1)$, the constant and
$\Nm(u)$ terms cancel.  The sum of the pairwise products of
$u,u^q,u^{q^2}$ equals $\Tr(u^{-1})$.  Consequently
\[
 \Nm(u+1)=\Tr(u)+\Tr(u^{-1})=\Tr(a),
\]
and hence
\begin{equation}\label{eq:cubic-law}
 \boxed{\ \Nm(a)=\Tr(a)^2.\ }
\end{equation}

Put $b=v+v^{-1}=a+1$ and $T=\Tr(a)$, and let
\[
 E_2=aa^q+aa^{q^2}+a^qa^{q^2}.
\]
The same identity for $b$, together with $\Tr(1)=1$, gives
$\Nm(b)=(T+1)^2=T^2+1$.  On the other hand,
\[
 \Nm(a+1)=\Nm(a)+E_2+T+1.
\]
Using \eqref{eq:cubic-law} therefore gives $E_2=T$.
The polynomial having $a,a^q,a^{q^2}$ as its roots is thus
\[
 Z^3+TZ^2+TZ+T^2=(Z+T)(Z^2+T).
\]
Every root of the right side belongs to $\F_q$, because squaring is a
bijection on $\F_q$.  In particular $a\in\F_q$.
Now $T=a$ and $\Nm(a)=a^3$, so \eqref{eq:cubic-law} becomes
$a^3=a^2$.  Thus $a=0$ or $a=1$.
The equation $u+u^{-1}=0$ forces $u=1$, and
$u+u^{-1}=1$ forces $u$ to have order $3$.
This proves the classification.

The norm-one group is $\mu_{q^2+q+1}$.  A common root for the bounded
grid would require both $u,v\ne1$, which the classification excludes.
Apply Lemma~\ref{lem:roots} and \eqref{eq:half-nullity}.
\end{proof}

\begin{corollary}\label{cor:cubic-primes}
If $p>3$ is prime and $3\mid\ord_p(2)$, then
\[
 d(p^k-1)=0\qquad\text{for every }k\geq1,
\]
whether or not $p$ is Wieferich.
\end{corollary}

\begin{proof}
Fix $k$ and write $f_k=\ord_{p^k}(2)=f_1p^e$, where
$f_1=\ord_p(2)$ and $e\geq0$.
Set $q=2^{f_k/3}$.
Then $q^3\equiv1\pmod{p^k}$.
Since $p\ne3$ and $p\nmid f_1$, the order of $q$ modulo $p$ is
exactly $3$: equivalently,
$\gcd(f_1,f_k/3)=f_1/3$.
Thus $q-1$ is a unit modulo $p^k$, and
\[
 p^k\mid q^2+q+1.
\]
Every $p^k$th root of unity lies in the cubic norm-one group.
The group $\mu_{p^k}$ contains no element of order $3$, so
Theorem~\ref{thm:cubic} excludes every relation
\eqref{eq:relation}.  The nullity is zero.
\end{proof}

\subsection{A quadratic trace--norm law from integer factorization}

\begin{theorem}[Two-factor exclusion]\label{thm:two-factors}
Let $r\geq1$ be odd, and put
\[
 q=2^r,\qquad s=2^{(r+1)/2},\qquad M_\pm=q+1\pm s.
\]
Suppose each of $u,v$ belongs to
$\mu_{M_-}\cup\mu_{M_+}$; the two coordinates need not belong to
the same subgroup.
Then \eqref{eq:relation} holds if and only if $u$ has order $5$
and $v=u^2$ or $u^{-2}$.
\end{theorem}

\begin{proof}
The integer identities
\begin{equation}\label{eq:aurifeuillian}
 M_-M_+=q^2+1,\qquad \gcd(M_-,M_+)=1
\end{equation}
follow from $s^2=2q$: the factors are odd and their difference $2s$
is a power of two.
In particular, $u^{q^2}=u^{-1}$, so $a=u+u^{-1}$ belongs to
$\F_{q^2}$.  Write
\[
 T(a)=a+a^q,\qquad N(a)=aa^q.
\]
Membership in either subgroup gives $u^{q+1}=u^s$ or $u^{-s}$.
Since $s$ is a power of two,
\[
 u^{q+1}+u^{-(q+1)}=a^s.
\]
Taking $q$th powers and using $u^{q^2}=u^{-1}$ gives
\[
 u^{q-1}+u^{-(q-1)}=(a^q)^s.
\]
Expanding $aa^q$ and adding these two identities proves
\begin{equation}\label{eq:quadratic-law}
 \boxed{\ N(a)=T(a)^s.\ }
\end{equation}
This proof works for either sign in $M_\pm$.

The element $b=v+v^{-1}$ satisfies the same law.  If $b=a+1$, then
$T(b)=T(a)$, so $N(b)=N(a)$.  But
\[
 N(a+1)=N(a)+T(a)+1.
\]
Hence $T(a)=1$ and \eqref{eq:quadratic-law} gives $N(a)=1$.
It follows that
\[
 a^q=a+1,\qquad a(a+1)=1,\qquad a^2+a+1=0.
\]
Substitute $a=u+u^{-1}$ and multiply by $u^2$ to obtain
\[
 u^4+u^3+u^2+u+1=0.
\]
Thus $u$ has exact order $5$.  Moreover $b=a+1=a^2$, so the
two roots of $Z^2+bZ+1$ are $u^2,u^{-2}$, giving the asserted
possibilities for $v$.
Conversely, a primitive fifth root and either of these choices satisfy
\eqref{eq:relation}.
\end{proof}

\begin{corollary}\label{cor:quadratic-primes}
Let $p$ be an odd prime with $v_2(\ord_p(2))=2$.
Then, for every $k\geq1$,
\[
 d(p^k-1)=
 \begin{cases}
  4,&p=5,\\
  0,&p\ne5.
 \end{cases}
\]
In particular, every prime $p\equiv5\pmod8$, other than $5$,
has zero nullity at all these levels, including any Wieferich primes
in that congruence class.
\end{corollary}

\begin{proof}
The quotient $\ord_{p^k}(2)/\ord_p(2)$ is a power of the odd
prime $p$.  Thus $\ord_{p^k}(2)=4r$ with $r$ odd.
The unique element of order two in $(\Z/p^k\Z)^\times$ is $-1$, so
\[
 2^{2r}\equiv-1\pmod{p^k}.
\]
With the notation of Theorem~\ref{thm:two-factors},
$p^k\mid M_-M_+$.  Since the two factors are coprime and $p^k$ is a
prime power, it divides one factor.
The theorem therefore applies to every pair in $\mu_{p^k}^2$.
For $p\ne5$ there are no solutions.
For $p=5$, only fifth roots occur, whose two trace values are the
roots of $X^2+X+1$; the nullity is $2\cdot2=4$.

Finally, if $p\equiv5\pmod8$, then $v_2(p-1)=2$ and Euler's
criterion gives $2^{(p-1)/2}\equiv-1\pmod p$.
Since $\ord_p(2)\mid p-1$, its two-adic valuation must be $2$.
\end{proof}

\subsection{Analytic proofs for both known Wieferich primes}

\begin{corollary}\label{cor:known-analytic}
Without either endpoint search,
\[
 \boxed{\ d(1093^k-1)=d(3511^k-1)=0
            \quad\text{for all }k\geq1.\ }
\]
\end{corollary}

\begin{proof}
For $1093$, use $1093\equiv5\pmod8$ and
Corollary~\ref{cor:quadratic-primes}.
For $3511$, use $\ord_{3511}(2)=1755$, which is divisible by $3$,
and Corollary~\ref{cor:cubic-primes}.
A compact certificate of the latter order is
\[
 2^{1755}\equiv1,\quad
 2^{585}\equiv756,\quad
 2^{351}\equiv1578,\quad
 2^{135}\equiv88\pmod{3511}.
\]
Since the prime divisors of $1755$ are $3,5,13$, these congruences
prove that the order is exactly $1755$.
\end{proof}

These are not arguments that presume $c_p=2$: both corollaries hold
at every level for any value of the Wieferich depth.
The factorization step in Corollary~\ref{cor:quadratic-primes} is also
where prime-power order matters.  One cannot put an arbitrary divisor
of $M_-M_+$ into one factor if it uses primes from both.

\section{An exact formula for the increment}\label{sec:increment}

The endpoint condition is usually expressed as the triviality of a gcd.
The following identity also measures the increase when that condition
fails, without omitting mixed collisions.

\begin{theorem}[Increment factorization]\label{thm:increment}
Let $M\mid N$ be positive odd integers.  Put
\[
 K(X)=\gcd\!\left(\frac{H_N(X)}{H_M(X)},H_N(X+1)\right),
 \qquad J(X)=\gcd(K(X),K(X+1)).
\]
All gcds are monic.  Then
\begin{equation}\label{eq:factor}
 G_N(X)=G_M(X)\frac{K(X)K(X+1)}{J(X)}
\end{equation}
and
\begin{equation}\label{eq:increment}
 \boxed{\ d(N-1)-d(M-1)=4\deg K-2\deg J.\ }
\end{equation}
In particular, $d(N-1)=d(M-1)$ if and only if $K=1$.
\end{theorem}

\begin{proof}
Write $A=T_N\setminus T_M$ and
\[
 C_N=T_N\cap(1+T_N),\qquad C_M=T_M\cap(1+T_M).
\]
By Lemma~\ref{lem:roots}, these are the root sets of $G_N$ and $G_M$.
The root set of $K$ is
\[
 S=\{\alpha\in A:\alpha+1\in T_N\}.
\]
A point of $C_N$ fails to belong to $C_M$ exactly when at least one
of $\alpha,\alpha+1$ is in $A$.  Therefore
\[
 C_N\setminus C_M=S\cup(1+S).
\]
The two sets on the right intersect in the root set of $J$.  All the
polynomials involved are square-free, so their union is represented by
$\operatorname{lcm}(K(X),K(X+1))=K(X)K(X+1)/J(X)$.
This union is disjoint from $C_M$, proving \eqref{eq:factor}.
Taking degrees and using \eqref{eq:half-nullity} proves
\eqref{eq:increment}.
Finally $\deg J\leq\deg K$, so a nonconstant $K$ gives a strictly
positive increment.
\end{proof}

More explicitly, put
\[
 a=\deg J,\qquad b=\deg K-\deg J.
\]
There are $a$ common roots for which both coordinates
$\alpha,\alpha+1$ are new, and $b$ common roots with the first
coordinate new and the second old.  Translation supplies another $b$
roots with the opposite orientation.  Hence the increment is
\[
 2a+4b.
\]
For a prime-power endpoint, ``new'' means order above $p$, not
necessarily one particular order when $c_p>2$.
Checking only new--new collisions cannot establish stabilization.

\begin{remark}[A genuine mixed example, not a prime-power counterexample]
Direct polynomial arithmetic for $M=17$ and $N=255$ gives
\[
 d(16)=8,\quad d(254)=140,\quad \deg K=62,\quad \deg J=58.
\]
Thus $a=58$, $b=4$, and the increment is
$2\cdot58+4\cdot4=132$.
An internal-collision calculation alone would miss the contribution
$16$ from the mixed pairs.  Since $255=3\cdot5\cdot17$, this example
does not contradict the prime-power conjecture.

One explicit mixed pair can be checked in
$\F_2[t]/(t^8+t^4+t^3+t^2+1)$, where $t$ has order $255$:
\[
 u=t^{17},\quad v=t^{120},\quad \ord(u)=15,\quad\ord(v)=17,
\]
\[
 u+u^{-1}=t^7+t^4+t+1,\qquad v+v^{-1}=t^7+t^4+t.
\]
Thus the two traces differ by one, even though only the second
coordinate has order dividing $17$.
\end{remark}

\section{Exactly what is already proved}\label{sec:descent}

For an odd prime $p$, write
\[
 f=\ord_p(2),\qquad c=v_p(2^f-1),\qquad
 \rho_i=\min\{r\geq1:2^r\equiv1\text{ or }-1\pmod{p^i}\}.
\]
The order-lifting formula and cyclicity of
$(\Z/p^i\Z)^\times$ give
\begin{equation}\label{eq:orders}
 \ord_{p^i}(2)=f p^{\max(0,i-c)},\qquad
 \rho_i=\rho_1p^{\max(0,i-c)}.
\end{equation}
Indeed, a multiple $ft$ of $f$ satisfies
$v_p(2^{ft}-1)=c+v_p(t)$.
The signed order is the ordinary order divided by two when the
ordinary order is even, and is the ordinary order otherwise.

If $u$ has exact order $p^i$, then
\begin{equation}\label{eq:trace-degree}
 [\F_2(u+u^{-1}):\F_2]=\rho_i.
\end{equation}
To see this, Frobenius to the power $r$ fixes $u+u^{-1}$ exactly when
$u^{2^r}$ is $u$ or $u^{-1}$.

\begin{theorem}[Known descent above the plateau]\label{thm:descent}
Every common root of $F_{p^k}(X)$ and $F_{p^k}(X+1)$ already belongs
to $T_{p^{\min(k,c)}}$.  In particular,
\[
 d(p^k-1)=d(p^c-1)\qquad(k\geq c).
\]
\end{theorem}

\begin{proof}
Write a common root and its translate as
\[
 \alpha=u+u^{-1},\qquad \alpha+1=v+v^{-1},
\]
where $u,v\ne1$ have exact orders $p^i,p^\ell$, respectively.
Since $\F_2(\alpha)=\F_2(\alpha+1)$,
\eqref{eq:orders}--\eqref{eq:trace-degree} imply either
$i,\ell\leq c$, or $i=\ell>c$.

Suppose $i=\ell>c$.  Set $L=p^i$ and $m=p^{i-c}$.
Taking $u$ as a primitive $L$th root, write $v=u^s$ with $p\nmid s$
and $1\leq s<L$.  Then $u$ is a zero of
\[
 P(X)=1+X+X^{L-1}+X^s+X^{L-s}.
\]
Let $h$ be the minimal polynomial of $u^m$ over $\F_2$.
Its degree is $f$, while the degree of $u$ is $fm$.
Therefore $h(X^m)$ is the minimal polynomial of $u$, and divides $P$.

Decompose polynomials uniquely by their exponent classes modulo $m$:
\[
 P(X)=\sum_{r=0}^{m-1}X^rP_r(X^m).
\]
Divisibility by $h(X^m)$ implies $h\mid P_r$ for every $r$.
But none of $1,-1,s,-s$ is zero modulo $m$, so $P_0=1$.
This is impossible for the nonconstant polynomial $h$.
Thus both orders are at most $p^c$.
The root sets are nested and every root of $F_{p^k}$ has multiplicity
two, proving the stated equality of nullities.
\end{proof}

\begin{corollary}
The desired identity holds for every non-Wieferich odd prime.
It also holds for $p=2$.
\end{corollary}

\begin{proof}
For a non-Wieferich odd prime, $c=1$, so apply
Theorem~\ref{thm:descent}.
For $p=2$, $F_{2^k}(X)=X^{2^k-1}$ is coprime to its translate.
Thus $d(2^k-1)=d(1)=0$.
\end{proof}

\section{The remaining assertion, with both types of collision}

\begin{problem}[Wieferich endpoint]\label{prob:endpoint}
For every Wieferich prime $p$, with $c=v_p(2^{\ord_p(2)}-1)$, prove
\begin{equation}\label{eq:missing}
 K_p(X):=\gcd\!\left(
       \frac{H_{p^c}(X)}{H_p(X)},H_{p^c}(X+1)
       \right)=1.
\end{equation}
\end{problem}

By Theorems~\ref{thm:increment} and~\ref{thm:descent}, this assertion
is equivalent to the original all-prime identity.
It is not a consequence of those theorems.
In fact, they give the exact possible discrepancy
\[
 d(p^c-1)-d(p-1)
 =4\deg K_p-2\deg\gcd(K_p(X),K_p(X+1)).
\]
By Corollaries~\ref{cor:cubic-primes} and
\ref{cor:quadratic-primes}, the unresolved primes can now be
restricted further to
\begin{equation}\label{eq:remaining-primes}
 c_p\geq2,\qquad 3\nmid\ord_p(2),\qquad
 v_2(\ord_p(2))\ne2.
\end{equation}
No claim that this set is empty is made.

For $p>3$, an equivalent finite-field formulation is
\begin{equation}\label{eq:endpoint}
 u,v\in\mu_{p^c},\quad
 1+u+u^{-1}+v+v^{-1}=0
 \quad\Longrightarrow\quad u,v\in\mu_p.
\end{equation}
Neither coordinate can equal $1$: that would force the other to have
order $3$.
Thus \eqref{eq:endpoint} is precisely the common-root condition of
Lemma~\ref{lem:roots}.
Notice that the hypothesis allows different orders for $u$ and $v$.

\subsection{Why the established proof does not extend through the plateau}

For $1\leq i\leq c$, all groups $\mu_{p^i}$ lie in the same field
$\F_{2^f}$, and their primitive traces all have degree $\rho_1$.
Equal degrees no longer imply equal levels.
In the residue-class proof, the proposed substitution has $m=1$,
so its decisive class $P_0$ is the whole polynomial, not the constant
$1$.

There is also no field-theoretic descent obtained simply by taking
$p$th powers.  On $\mu_{p^i}$ this operation has a kernel of order $p$;
it is not a field automorphism.  The fact that $u^p,v^p$ have smaller
orders does not imply that they satisfy the same additive equation.
Likewise, a norm identity does not make all conjugates of an element
vanish at a chosen prime above two.  On the plateau the cyclotomic
extension splits at two, rather than supplying the nontrivial residue
field extension used by the off-plateau proof.

\subsection{What a counterexample certificate must contain}

A counterexample can be presented without constructing a large grid.
Choose a prime $p>3$, an integer $2\leq j\leq c_p$, an irreducible
$h\in\F_2[T]$, and an integer $b$ with $1\leq b<p^j$.  In
$\F_2[T]/(h)$, let $t$ be the residue of $T$ and verify
\[
 t^{p^j}=1,\qquad t^{p^{j-1}}\ne1,\qquad
 1+t+t^{p^j-1}+t^b+t^{p^j-b}=0.
\]
These equalities exhibit a new common root and imply
$d(p^j-1)>d(p-1)$.
Conversely, from any counterexample choose a coordinate of largest
order as $t$; this produces such a certificate.
One must not require $p\nmid b$, since a mixed-order counterexample
could have $p\mid b$.

\section{Exact computations and their scope}\label{sec:computations}

The existing search in \cite{code} uses characteristic-two field
arithmetic, not the conjectured prime-power reduction.
Its 64-bit fingerprint is an additive projection: an exact equality
necessarily survives projection.  Thus an empty fingerprint-candidate
set is a rigorous absence certificate, not a probabilistic conclusion.
Any surviving candidates are checked in the full field.

Both known Wieferich endpoint searches were rerun in this investigation:
\[
\begin{array}{c|rrrr}
 p&c_p&\ord_p(2)&\text{fingerprint candidates}&
       d(p^2-1)=d(p-1)\\ \hline
 1093&2&364&0&0\\
 3511&2&1755&0&0
\end{array}
\]
The depth $c_p=2$ is verified by
$2^{\ord_p(2)}\equiv1\pmod{p^2}$ but not modulo $p^3$.
The zero endpoint nullities imply zero base nullities by root-set
inclusion.  Reproducible commands from the repository root are
\begin{verbatim}
python code\ai_prototypes\five_term_char2.py 1093 2
python code\ai_prototypes\five_term_char2.py 3511 2
\end{verbatim}
These are now independent checks of Corollary~\ref{cor:known-analytic},
not prerequisites for it.

The trace--norm laws were also checked directly in finite fields.
For the cubic law, $q=2,4,8,16,32$ gives only the classified degenerate
relations: respectively $0,4,0,4,0$ ordered pairs.
For the two-factor law, $r=1,3,5,7$ gives
$(M_-,M_+)=(1,5),(5,13),(25,41),(113,145)$.
In each union of the two subgroups there are exactly eight ordered
relations, all from the fifth roots as proved above.

For small examples, direct Fibonacci gcd calculations give
\[
\begin{array}{c|rrr}
 p&d(p-1)&d(p^2-1)&d(p^3-1)\\ \hline
 3&0&0&0\\
 5&4&4&4\\
 7&0&0&0\\
 11&0&0&0\\
 17&8&8&8\\
 31&20&20&20
\end{array}
\]
These calculations use \texttt{\_adjacent\_fibonacci\_pair} and
\texttt{GF2Polynomial.gcd}, not the optimized
\texttt{grid\_nullity} shortcut for proven prime-power cases.
The factorization in Theorem~\ref{thm:increment} was also checked
directly for every proper odd divisor pair $M\mid N$, $3\leq N\leq301$:
398 pairs, including 13 with a nonzero mixed contribution.
The proof of the theorem is the root-set argument, not this finite
check.

\section{Conclusion}\label{sec:conclusion}

Two arithmetic classes are now covered without a
non-Wieferich hypothesis:
$3\mid\ord_p(2)$ and $v_2(\ord_p(2))=2$.
In particular, both known Wieferich primes have analytic all-level
proofs, rather than only computational endpoint proofs.
The increment factorization additionally measures both kinds of
possible new collision.

The all-prime question nevertheless remains unresolved by this
investigation.  The missing step is \eqref{eq:missing}, equivalently
\eqref{eq:endpoint}, for a Wieferich prime satisfying
\eqref{eq:remaining-primes}.
Neither trace--norm theorem addresses all of those primes, and no
counterexample certificate of the required kind has been found.

Further attempts at conductor descent, rather than additional
zero-nullity families, are recorded in
\path|tex\prime_power_endpoint\conductor_descent.tex|.
That addendum gives an obstruction to direct powering in the
regular-map formulation and a single-candidate-phase criterion for
mixed collisions when $H_p$ is irreducible; it does not claim a
general proof.

\begin{thebibliography}{9}

\bibitem{finite-fields}
W.~Boyles, \emph{Fibonacci Polynomials \& Lights Out}.
Repository manuscript:
\path|tex\finite_fields\finite_fields.tex|.
See ``Orders of Prime Powers'' and ``Applications to $d(n)$''.

\bibitem{plateau}
W.~Boyles, \emph{Reciprocal Five-Term Relations among $p^j$-th Roots
of Unity in Characteristic Two: Reductions, Obstructions, and the
Wieferich Plateau}.
Repository manuscript:
\path|tex\five_term_plateau\five_term_plateau.tex|.
See ``The single endpoint statement'', ``Two collision norms'',
and ``Computations''.

\bibitem{code}
Repository implementations:
\path|code\kernel_size.py|,
\path|code\polynomials.py|,
\path|code\ai_prototypes\five_term_char2.py|,
and \path|code\ai_prototypes\wieferich_check.py|.

\end{thebibliography}
\end{document}
